% Abstract page (follows the cover page).
% Title is reused from coverpage.tex (\coverTitle).

\clearpage
\thispagestyle{empty}
\begin{center}
  \vspace*{0.5in}
  {\Large\bfseries\coverTitle\par}
  \vspace{0.25in}
  {\large Levi Le\par}
  \vspace{0.4in}
  {\bfseries Abstract\par}
\end{center}

\vspace{0.1in}
\noindent
Mobile robots usually localize with a dedicated range sensor such as a LiDAR, which adds cost,
weight, and power draw that are hard to justify on a small, battery-powered platform. This thesis
asked whether monocular depth estimation (MDE), which predicts the depth of every pixel from a
single camera image, is good enough to take the place of such a sensor. To answer this question,
we built the Cyclops rover, a small wheeled robot whose only view of its surroundings is a single
front-mounted camera. Each camera frame was passed to Depth Anything V2, a vision transformer
model that predicts relative depth. Because relative depth has an unknown scale and shift that
change from frame to frame, we recovered metric depth on every frame by fitting the scale and
shift to points on the floor whose depth is fixed by the camera's mounting. The metric depth map
was then back-projected into a point cloud, stripped of the floor, and reduced to a planar range
scan of the same form as a LiDAR scan. This scan served as the measurement for Monte Carlo
localization in a known occupancy grid map, alongside wheel odometry, and the resulting pose
estimate guided A* path planning toward a goal. In tests inside a $1.105$~m $\times$ $1.82$~m
enclosure, the rover reached its goal in 10 of $10$ trials
with localization that closely matches it real pose. These results show
that, on a flat, well-lit floor and at short range, a single camera can replace a dedicated range
sensor for localizing a small robot.

\vspace{0.2in}
\noindent
\textit{Keywords}: Monocular depth estimation, Monte Carlo localization, Vision transformer,
Depth Anything V2, Mobile robot

\clearpage

%%% Local Variables: ***
%%% mode: latex ***
%%% TeX-master: "thesis.tex" ***
%%% End: ***
